% ---------------------------------------------------------------
\section{Conclusion}\label{sec:conclusion}
% ---------------------------------------------------------------

The Bailout Protocol contributes a single, architecturally precise idea: in any multi-agent system built on the \bxthree{} Framework, every exception condition that exceeds machine-actor provisioned authority must — by architectural compulsion, not by machine-actor choice — reach a named human principal who can bind it. This is not an error-handling strategy selected at runtime. It is a structural constraint on the architecture of delegation itself, embedded in the deterministic state machine of every \bxthree{} node and enforced by the Fact Layer's pre-commit hook, which admits no machine-actor override.

The contribution is operational and architectural, not merely philosophical. The protocol delivers four concrete guarantees that existing escalation hierarchies do not provide. First, it is \emph{mandatory} rather than optional: the state machine's transition to \texttt{ESCALATING} is triggered by the observational boundary condition (TC) and by the Fact Layer's timeout enforcement, not by a machine actor's willingness to surface a failure. Second, it is \emph{non-suppressible}: once in \texttt{ESCALATING} state, the bypass mechanism eliminates every intermediate machine actor from the communication path to $h(n)$, so no machine actor can absorb, redirect, or delay the escalation. Third, it is \emph{observable}: every protocol event — trigger evaluations, state transitions, pathway selections, acknowledgments, and timeout conditions — is committed to the append-only Forensic Ledger with SHA-256 hash-chain integrity, making retrospective verification of the protocol's correct execution deterministic rather than inferential. Fourth, it is \emph{live}: the timeout architecture ($T_{\bounds}$, $T_{\prop}$, $T_{\cap}$, $T_{\human}$) defines a finite upper bound $T_{\max}$ on wall-clock time from trigger firing to human acknowledgment, subject only to the human anchor's availability and the Pathway Health Registry's pathway quality.

The protocol does not stand alone. It is the fifth and final pillar of the \bxthree{} Framework's upstream accountability architecture. Loop Isolation, Recursive Spawning, Spatial Firewall, and Sandbox Gate address the exception conditions anticipated at design time — those within the machine-only execution graph and resolvable by bounded agentic reasoning. The Bailout Protocol addresses the complementary remainder: conditions not anticipated, that exceed the operating range of any machine actor, and that cannot be resolved within the machine-only execution graph. Together, the five pillars constitute structurally complete coverage of the exception space: every condition that can arise in a \bxthree{} node either is resolved within the machine-only execution graph by the appropriate preceding pillar, or is compelled by the Bailout Protocol to reach a human who can bind it. The upstream accountability guarantee — that \bxthree{} systems fail upward into human consciousness, never downward into autonomous chaos — is unconditional precisely because this coverage is complete.

The \bxthree{} Framework's role in this contribution is structural: it provides the three-layer accountability architecture — \purpose{} carrying intent and final authority, \bounds{} carrying bounded execution, and \fact{} carrying deterministic enforcement — without which the Bailout Protocol's guarantees cannot be realized. The immutable parent pointer, the non-adjustable confidence threshold, the bypass mechanism, and the Forensic Ledger are each products of the three-layer model. The protocol and the framework are not separable: the framework defines the structural properties that make the protocol possible; the protocol is the runtime mechanism that makes the framework's upstream accountability guarantee operative. One cannot be understood without the other.

The limitations enumerated in Section~\ref{sec:limitations} are real and acknowledged. Human response latency under burst loads, adversarial pathway integrity, scalability at high delegation depth, and compulsory escalation as a denial-of-service vector each represent a surface area that the current specification does not fully address. Future work — formal verification, adaptive timeout, distributed human anchor pools, simultaneous multi-pathway escalation, and cryptographic pathway attestation — is targeted precisely at these surfaces, and each extension is designed to preserve the bypass property and the mandatory escalation guarantee rather than relax them. The protocol is a foundation, not a finished edifice. Its structural core is sound and its extensions are well-defined. The foundation is what matters operationally today.

\endinput